\documentclass[10pt,a4paper]{report}
\usepackage[utf8]{inputenc}
\usepackage{amsmath}
\usepackage{amsfonts}
\usepackage{amssymb}
\usepackage{amsthm}
\usepackage{hyperref}

\usepackage{multicol}
\usepackage{fancyhdr}
\usepackage[inline]{enumitem}
\usepackage{tikz}
\usepackage{tikz-cd}
\usetikzlibrary{calc}
\usetikzlibrary{shapes.geometric}
\usepackage[margin=0.5in]{geometry}
\usepackage{xcolor}
\usepackage{ esint }

\hypersetup{
    colorlinks=true,
    linkcolor=blue,
    filecolor=magenta,      
    urlcolor=cyan,
    pdftitle={Tensors},
    pdfpagemode=FullScreen,
    }

%\urlstyle{same}

\newcommand{\CLASSNAME}{Math 5230 -- Partial Differential Equations}
\newcommand{\STUDENTNAME}{Paul Carmody}
\newcommand{\ASSIGNMENT}{Homework \#4 }
\newcommand{\DUEDATE}{October 29, 2025}
\newcommand{\SEMESTER}{Fall 2025}
\newcommand{\SCHEDULE}{MW 2:00 -- 3:15}
\newcommand{\ROOM}{Crawford 332}

\newcommand{\MMN}{M_{m\times n}}
\newcommand{\FF}{\mathcal{F}}

\pagestyle{fancy}
\fancyhf{}
\chead{ \fancyplain{}{\CLASSNAME} }
%\chead{ \fancyplain{}{\STUDENTNAME} }
\rhead{\thepage}
\input{../MyMacros}

\newcommand{\RED}[1]{\textcolor{red}{#1}}
\newcommand{\BLUE}[1]{\textcolor{blue}{#1}}

\newcommand{\CURL}{\operatorname{curl}}
\newcommand{\DIV}{\operatorname{div}}

\begin{document}

\begin{center}
	\Large{\CLASSNAME -- \SEMESTER} \\
	\large{ w/Professor Snelson}
\end{center}
\begin{center}
	\STUDENTNAME \\
	\ASSIGNMENT -- \DUEDATE\\
\end{center} 

\noindent Please show your work and justify all answers.

\begin{enumerate}

\item If $F$ and $G$ are any twice-differentiable functions, show that 
\begin{align*}
	u(x,t) = F(x+t)+G(x-t),
\end{align*}solves the one-dimensional wave equation $u_{tt}-u_{xx}=0$.

\BLUE{\begin{align*}
	u_t &= F'(x+t) -G'(x-t) & \quad u_{tt} &= F''(x+t) + G''(x-t)\\
	u_x &= F'(x+t) + G'(x-t) & \quad u_{xx} &= F''(x+t) + G''(x-t)
\end{align*}\begin{align*}
	u_{tt}-u_{xx} &= F''(x+t) + G''(x-t) - (F''(x+t) + G''(x-t)) \\
	&=0
\end{align*}
}

\item Let $P$ be a parallelogram in $(x,t)$ space whose sides have slope -1 and 1.  Let $p_1,p_2,p_3,p_4$ be the four vertices of $P$, listed in counterclockwise order (Note that each $p_i$ is a point in $(x,t)$ space.)  Show that, if $u$ is a solution of the wave equation and $P$ lies within the domain of $u$, then 
\begin{align*}
	u(p_1)+u(p_3)=u(p_2)+u(p_4)
\end{align*}(Hint: you are allowed to use the converse to problem 1: any solution to $u_{tt}-u_{xx}=0$ can be written as $F(x+t)+G(x-t)$ for some twice-differentiable functions $F$ and $G$.)

\BLUE{
\begin{align*}
	u(x,t) &= F(x+t) + G(x-t)
\end{align*}Let $\alpha = x+t$ and $\beta = x-t$. Then,
\begin{align*}
	u(\alpha, \beta) &= F(\alpha)+G(\beta)
\end{align*}Notice that for sides with slope equal to 1 we have 
\begin{align*}
	m &= 1  & \quad x_1+t_1 & =x_2+t_2 & \quad \alpha_1 &= \alpha_2\\
	 & \quad  & x_3+t_3 &= x_4+t_4 & \quad  \alpha_3&=\alpha_4 \\
	m &= -1  & \quad x_1-t_1 & =x_4-t_4 & \quad \beta_1 &= \beta_4\\
	 & \quad  & x_2-t_2 &= x_3-t_3 & \quad  \beta_2&=\beta_3 
\end{align*}Thus,
\begin{align*}
	u(p_1) &= F(\alpha_1)+G(\beta_1), & \quad u(p_2) &= F(\alpha_1)+G(\beta_2) \\
	u(p_3) &= F(\alpha_3)+G(\beta_2), & \quad u(p_2) &= F(\alpha_3)+G(\beta_1) \\
\end{align*}Then we can see that
\begin{align*}
	u(p_1)+u(p_3) &= F(\alpha_1)+G(\beta_1) + F(\alpha_3)+G(\beta_2) \\
	u(p_2)+ u(p_4) &= F(\alpha_1)+G(\beta_2) + F(\alpha_3)+G(\beta_1)
\end{align*}which are equal.
}

\item Consider the wave equation on a  half-infinite string:
\begin{align*}
	\BINDEF{u_{tt}-u_{xx}=0, & t>0,x>0, \\u(x,0)=0}{u_t(x,0)=0\\u(0,t)=h(t).}
\end{align*}(In class, we looked at the case where the boundary condition at $x=0$ was zero, and the initial conditions were nonzero.  This is the opposite case.  Odd reflection will not work here because the boundary condition isn't zero.)

Derive a solution formula for this problem.  (Hint: consider the cases $x>t$ and $x\le t$ separately.  In the case $x\le t$, use the result of problem 2.)

\BLUE{
Applying the initial conditions and recall that $	u(x,t) = F(x+t) + G(x-t) = 0$
\begin{align*}
	F(x+t) &= -G(x-t) \\ 
	u(x,0) &= F(x)+G(x) = 0 \implies F(x) = - G(x) & (*)\\
	u(0,t) &= F(t)+G(-t) = h(t) \\
	u(0,t) &= F(-t)+G(t) = h(t) \text{  for } x< t
\end{align*}Taking the first partial with respect to $t$
\begin{align*}
	u_t(x,t) &= F'(x+t)-G'(x-t)\\
	u_t(x,0) &= F'(x)-G(x)=0 \implies F'(x) = G'(x)
\end{align*}but from $(*)$ above we have $F'(x) = -G'(x) = -F'(x)$.  Therefore $F'(x)=0$ and $F(x)$ is a constant.  That leaves $G(x-t)$ which only makes sense when $x\ge t$.  Hence,
\begin{align*}
	u(x,t) &= \BINDEF{h(x,t) & x \ge t}{0 & x <0}
\end{align*}
}

\item Recall d'Alembert's formula for the one-dimensional wave equation:
\begin{align*}
	u(x,t) = \frac{1}{2}[g(x+t)+g(x-t)] + \frac{1}{2}\int_{x-t}^{x+t} h(y)\, dy.
\end{align*}Find ``fundamental solutions $G_0(x,t)$ and $G_1(x,t)$ so that this formal can be written convolution form:
\begin{align*}
	u(x,t) = \int_\R G_0(x-y,t)g(y)\,dy+\int_\R G_1(x-y,t)h(y)\,dy
\end{align*}(Hint:  one of these will involve delta functions Recall the definition property of the delta function: $\int_\R\delta(y-c)f(y)\,dy=f(c)$ for any function $f$.)

\BLUE{
We start by finding convolving for the first term.
\begin{align*}
	\frac{1}{2}[g(x+t)+g(x-t)] &= \int_\R [\delta(y-(x+t))+\delta(y+(x-t))]g(y)\,dy \\
	&= \int_\R \frac{1}{2}[\delta(x-y+t)+\delta(x-y-t)]g(y)\,dy \\
	G_0 (x,t) &= \frac{1}{2}[\delta(x+t)+\delta(x-t)]
\end{align*}
}

\newcommand{\EE}{\textbf{E}}
\newcommand{\BB}{\textbf{B}}
\item Evans Chapter 2, Problem 21(a).  Assume $\EE=(E^1,E^2,E^3)$ and $\BB= (B^1, B^2, B^3)$ solve Maxwell's equations
\begin{align*}
	\BINDEF{\EE_t = \CURL\BB, & \BB_t = -\CURL \EE}{\DIV \BB = \DIV \EE = 0.}
\end{align*}Show \begin{align*}
	\EE_{tt}-\Delta \EE = 0 & \quad \BB_{tt} - \Delta \BB = 0.
\end{align*}

\BLUE{
Let us first take the derivative with respect to time of the first two identities.
\begin{align*}
	\EE_{tt} &=	 \CURL \BB_t &\quad \BB_{tt} &= - \CURL \EE_t \\
	&= \nabla \times (-\nabla \times \EE) &\quad &= - \nabla \times (\nabla \times \BB) \\
	&= -\nabla(\DIV \EE)+\Delta \EE &\quad &= -\nabla(\DIV \BB) + \Delta \EE \\
	&= \Delta \EE &\quad &= \Delta \BB
\end{align*}
}

\end{enumerate}

\end{document}
